\documentclass[10pt,a4paper]{amsart}
\usepackage[margin=19mm]{geometry}
\usepackage{amsmath,amssymb,mathtools}
\usepackage[scr=rsfs,cal=euler]{mathalfa}
\usepackage{xfrac}
\usepackage[unicode,hidelinks,bookmarks=false]{hyperref}
\newcommand{\ie}{i.e.\ }
\newcommand{\cf}{cf.\ }
\newcommand{\resp}{respectively\ }
\newcommand{\Subsection}{\subsection}
\providecommand{\nobreakdash}{\nobreak\hbox{-}}
\setlength{\emergencystretch}{2em}
\multlinegap0pt
\usepackage{amssymb}
\usepackage{enumerate}
\usepackage[shortlabels]{enumitem}
\usepackage{tikz}
\newtheorem{lem}{Lemma}
\title{The logarithmic p-Laplacian on hyperbolic spaces}
\author{Jorge J. Betancor, Lourdes Rodr\'{\i}guez-Mesa}
\date{Source: arXiv:2608.01079v1 (CC BY 4.0)}
\begin{document}
\maketitle

\noindent\emph{Source and licence.} The logarithmic p-Laplacian on hyperbolic spaces. Version arXiv:2608.01079v1 (CC BY 4.0), distributed by arXiv under the Creative Commons Attribution 4.0 International licence, http://creativecommons.org/licenses/by/4.0/, as recorded in the arXiv licence metadata for this exact version and shown on the abstract page https://arxiv.org/abs/2608.01079v1. Full licence notice: \texttt{licenses/arxiv-2608.01079-CC-BY-4.0.txt}.\par\smallskip

\noindent\textit{Editorial scope.} Counted unit: the article's lem key (source0.tex lines 332--363). The article's complete original proof of it is delivered with it (source0.tex lines 364--514, one \texttt{proof} environment of 151 lines). No mathematics is written, repaired or re-derived: the statement and the proof are the article's own lines, in the article's own notation, and no retained line is substituted or re-flowed.  No further source-local context is needed: every cross-reference of every retained line resolves inside the two delivered spans.  Nothing was omitted from any retained span, so the package carries no omission record.  No retained line contains a citation, so the wrapper carries no bibliography and no bibliography entry.  No retained line contains a figure environment, an image-inclusion command or drawing code, so no image is delivered and the package declares no asset dependency.  Editorial formatting only, in addition: an AMS \texttt{amsart} preamble that restates the article's own declarations and macros used by the retained lines, namely the environments \texttt{lem} on separate counters, together with the standard packages those lines need.  The retained environments print in this document as Lemma 1 in order of appearance, whereas the article prints the same items with the numbering of its own full document.  The complete native source of the article as distributed in the arXiv e-print for this exact version is bundled unchanged as \texttt{sources/206594/source0.tex}.
\begin{lem}\label{key}
For every $m\in \mathbb N$ we have 
\begin{equation}\label{formderiv}
\Big(\frac{1}{\sinh \rho }\frac{\partial }{\partial\rho }\Big)^m=\sum_{k=1}^m\frac{T_{m,k}(\rho )}{(\sinh \rho)^{2m-k}}\Big(\frac{1}{\rho }\frac{\partial }{\partial\rho }\Big)^k,\quad \rho >0,
\end{equation}
where, for each $k=1,\ldots, m$, 
\begin{equation}\label{Tmk}
T_{m,k}(\rho )=\sum_{j=0}^k Q_j^{m,k}(\rho )\rho ^j,\quad \rho >0,
\end{equation}
with $Q_j^{m,k}(\rho )=R_j^{m,k}(\sinh \rho ,\cosh \rho)$, $\rho >0$, where $R_j^{m,k}$ is a homogeneous polynomial of degree $m-k$, for $j=0,\ldots , k$. In particular, $T_{m,m}(\rho )=\rho  ^m$, $\rho >0$.

Moreover, for every $m\in \mathbb N$, and $k=1,\ldots, m$, the following properties hold:

(i) For certain $(b_i^{m,k})_{i=0}^\infty\subseteq \mathbb R$, 
$$
T_{m,k}(\rho )=\sum_{i=0}^\infty b_i^{m,k}\rho ^{2m-k+2i},\quad \rho >0,
$$
and the series is absolutely convergent for each $\rho >0$. Note that if $k\in \mathbb N$ is even (respectively, odd) $T_{m,k}$ can be extended as an even function (respectively, odd function).

(ii) For each $j=0,\ldots,k$, there exists $c_j^{m,k}\in \mathbb R$ such that
$$
\Big|\frac{Q_j^{m,k}(\rho )}{(\sinh \rho )^{2m-k}}-c_j^{m,k}e^{-m\rho }\Big|\leq Ce^{-(m+2)\rho},\quad \rho \geq 1,
$$

where $C$ does not depend on $j$ or $k$. Furthermore, $|c_k^{m,k}|=(-1)^{m+k}c_k^{m,k}$.

(iii) For each $j=0,\ldots,k$, 
$$
\Big|\frac{Q_j^{m,k}(\rho )}{(\sinh \rho )^{2m-k}}\Big|\leq Ce^{-m\rho },\quad \rho \geq 1,
$$
for certain $C>0$ that does not depend on $k$ or $j$.
\end{lem}

\begin{proof}
We proceed by induction on $m\in \mathbb N$ to establish the formula \eqref{formderiv}. When $m=1$, it is clear that property holds by taking $T_{1,1}(\rho )=\rho $, $\rho >0$.
Suppose that \eqref{formderiv} is satisfied for a certain $m\in \mathbb N$ and that $T_{m,k}$, $k=1,\ldots, m$, is given by \eqref{Tmk}.

Then, we can write
\begin{align*}
 \Big(\frac{1}{\sinh \rho }\frac{\partial}{\partial \rho}\Big)^{m+1}&=\sum_{k=1}^m\left[\frac{\rho T_{m,k}(\rho )}{(\sinh \rho)^{2m+1-k}}\Big(\frac{1}{\rho }\frac{\partial }{\partial \rho} \Big)^{k+1}+ \frac{T_{m,k}'(\rho )}{(\sinh \rho)^{2m+1-k}}\Big(\frac{1}{\rho }\frac{\partial}{\partial \rho }\Big)^k\right.\\
 &\quad \left.-(2m-k)\frac{T_{m,k}(\rho )\cosh \rho }{(\sinh \rho)^{2m+2-k}}\Big(\frac{1}{\rho }\frac{\partial }{\partial \rho}\Big)^k\right]\\
 &=\sum_{k=2}^{m+1}\frac{\rho T_{m,k-1}(\rho )}{(\sinh \rho)^{2m+2-k}}\Big(\frac{1}{\rho }\frac{\partial }{\partial \rho}\Big)^k\\
 &\quad + \sum_{k=1}^m\frac{T_{m,k}'(\rho )\sinh \rho -(2m-k)T_{m,k}(\rho )\cosh \rho }{(\sinh \rho)^{2m+2-k}}\Big(\frac{1}{\rho}\frac{\partial}{\partial \rho}\Big)^k\\
 &=\sum_{k=1}^{m+1}\frac{T_{m+1,k}(\rho )}{(\sinh \rho )^{2(m+1)-k}}\Big(\frac{1}{\rho}\frac{\partial}{\partial \rho}\Big)^k,
\end{align*}
where
\begin{align}\label{T1}
&T_{m+1,1}(\rho )=T_{m,1}'(\rho )\sinh \rho -(2m-1)T_{m,1}(\rho )\cosh \rho,\quad \rho >0,\\
&T_{m+1,m+1}(\rho )=\rho T_{m,m}(\rho)=\rho ^{m+1},\quad \rho >0,\label{Tm}
\end{align}
and, for $k=2,\ldots, m$, 
\begin{equation}\label{Tk}
T_{m+1,k}(\rho )=\rho T_{m,k-1}(\rho )+T_{m,k}'(\rho )\sinh \rho -(2m-k)T_{m,k}(\rho )\cosh \rho,\quad \rho >0.
\end{equation}
Let $k=1,\ldots,m$. From \eqref{Tmk} we get
$$
T_{m,k}'(\rho )=(Q_k^{m,k})'(\rho )\rho ^k+\sum_{j=0}^{k-1}\big((Q_j^{m,k})'(\rho )+(j+1)Q_{j+1}^{m,k}(\rho )\big)\rho ^j,\quad \rho >0.
$$
Then, $T_{m+1,k}(\rho)=\sum_{j=0}^kQ_j^{m+1,k}(\rho )\rho ^j$, $\rho >0$, $k=1,\ldots, m+1$, where, for each $\rho >0$,
\begin{align}\label{Q11}
Q_j^{m+1,m+1}(\rho)&=\left\{
\begin{array}{ll}
0,&j=0,\ldots, m,\\[0.2cm]
1,&j=m+1,
\end{array}
\right.\nonumber\\
 Q_j^{m+1,1}(\rho)&=\left\{
\begin{array}{ll}
((Q_0^{m,1})'(\rho)+Q_1^{m,1}(\rho ))\sinh \rho-(2m-1)Q_0^{m,1}(\rho )\cosh \rho,&j=0,\\[0.2cm]
(Q_1^{m,1})'(\rho )\sinh \rho-(2m-1)Q_1^{m,1}(\rho )\cosh \rho,&j=1,
\end{array}
\right.
\end{align}
and, for $k=2,\ldots, m$,
\begin{equation}\label{Qkk}
Q_j^{m+1,k}(\rho)=\left\{
\begin{array}{ll}
((Q_0^{m,k})'(\rho)+Q_1^{m,k}(\rho ))\sinh \rho-(2m-k)Q_0^{m,k}(\rho )\cosh \rho,&j=0,\\[0.2cm]
Q_{k-1}^{m,k-1}(\rho )+(Q_k^{m,k})'(\rho )\sinh \rho-(2m-k)Q_k^{m,k}(\rho )\cosh \rho,&j=k,
\end{array}
\right.
\end{equation}
and, if $j=1,\ldots, k-1$, 
$$
Q_j^{m+1,k}(\rho)=Q_{j-1}^{m,k-1}(\rho)+((Q_j^{m,k})'(\rho)+(j+1)Q_{j+1}^{m,k}(\rho ))\sinh \rho-(2m-k)Q_j^{m,k}(\rho )\cosh \rho.
$$

Since for each $j=0,\ldots, k$, we have that $Q_j^{m,k}(\rho)=R_j^{m,k}(\sinh \rho ,\cosh \rho )$, $\rho >0$, where $R_j^{m,k}$ is a homogeneous polynomial of degreee $m-k$, and, for every $r=0,\ldots, m-k$, and $\rho >0$,
$$
\frac{d}{d\rho }\big [(\sinh \rho )^r(\cosh \rho )^{m-k-r}\big]=r(\sinh \rho )^{r-1}(\cosh \rho )^{m+1-k-r}+(m-k-r)(\sinh \rho )^{r+1}(\cosh \rho )^{m-k-r-1},
$$
it is clear that $(Q_j^{m,k})'\sinh \rho$, $Q_j^{m,k}\sinh \rho$, and $Q_j^{m,k}\cosh \rho$, $j=0,\ldots, k$, can be represented by means of homogeneous polynomials of degree $m+1-k$. Also, $R_j^{m,k-1}$, $j=0,\ldots, k-1$, and $k=2,\ldots, m$, are homogeneous polynomials of degree $m-(k-1)=m+1-k$. By using these arguments and from \eqref{Q11} and \eqref{Qkk} we conclude that for $k=1,\ldots, m+1$, and $j=0,\ldots, k$, $Q_j^{m+1, k}(\rho )=R_j^{m+1,k}(\sinh \rho ,\cosh \rho )$, $\rho >0$, with $R_j^{m+1,k}$ a homogeneous polynomial of degree $m+1-k$. 

Let us next prove $(i)$. We again proceed by induction on $m\in \mathbb N$, and use \eqref{T1}-\eqref{Tk}. It is clear for $m=1$. Suppose that $(i)$ holds for certain $m\in \mathbb N$. Note that, since $T_{m+1,m+1}(\rho )=\rho ^{m+1}$, $\rho >0$, it follows that $b_0^{m+1,m+1}=1$, and $b_i^{m+1,m+1}=0$, $i\in \mathbb N$. On the other hand by \eqref{T1} we get
\begin{align*}
    T_{m+1,1}(\rho )&=\sinh \rho \sum_{i=0}^\infty b_i^{m,1}(2m-1+2i)\rho ^{2m-2+2i}-(2m-1)\cosh \rho \sum_{i=0}^\infty b_i^{m,1}\rho ^{2m-1+2i}\\
    &=\sum_{i,r=0}^\infty b_i^{m,1}\frac{2m-1+2i}{(2r+1)!}\rho ^{2m-1+2(i+r)}-\sum_{i,r=0}^\infty b_i^{m,1}\frac{2m-1}{(2r)!}\rho ^{2m-1+2(i+r)}\\
    &=(2m-1)b_0^{m,1}\sum_{r=1}^\infty \Big(\frac{1}{(2r+1)!}-\frac{1}{(2r)!}\Big)\rho ^{2m-1+2r}\\
    &\quad +\sum_{i=1}^\infty b_i^{m,1}\sum_{r=0}^\infty \Big(\frac{2m-1+2i}{(2r+1)!}-\frac{2m-1}{(2r)!}\Big)\rho ^{2m-1+2(i+r)}\\
    &=(2m-1)b_0^{m,1}\sum_{r=0}^\infty \Big(\frac{1}{(2r+3)!}-\frac{1}{(2r+2)!}\Big)\rho ^{2m+1+2r}\\
    &\quad +\sum_{i=0}^\infty b_{i+1}^{m,1}\sum_{r=0}^\infty \Big(\frac{2m+1+2i}{(2r+1)!}-\frac{2m-1}{(2r)!}\Big)\rho ^{2m+1+2(i+r)}\\
     &=(2m-1)b_0^{m,1}\sum_{r=0}^\infty \Big(\frac{1}{(2r+3)!}-\frac{1}{(2r+2)!}\Big)\rho ^{2m+1+2r}\\
    &\quad +\sum_{r=0}^\infty \sum_{i=r}^\infty b_{i-r+1}^{m,1}\Big(\frac{2m+1+2(i-r)}{(2r+1)!}-\frac{2m-1}{(2r)!}\Big)\rho ^{2m+1+2i}\\
     &=(2m-1)b_0^{m,1}\sum_{r=0}^\infty \Big(\frac{1}{(2r+3)!}-\frac{1}{(2r+2)!}\Big)\rho ^{2m+1+2r}\\
    &\quad +\sum_{i=0}^\infty \sum_{r=0}^i b_{i-r+1}^{m,1}\Big(\frac{2m+1+2(i-r)}{(2r+1)!}-\frac{2m-1}{(2r)!}\Big)\rho ^{2m+1+2i}\\
    &=\sum_{i=0}^\infty b_i^{m+1,1}\rho ^{2(m+1)-1+2i},\quad \rho >0,
\end{align*}
for certain $(b_i^{m+1,1})_{i=0}^\infty \subseteq \mathbb R$, being the series absolutely convergent for each $\rho >0$. Analogously, when $k=2,\ldots, m$, by using \eqref{Tk} we can write
\begin{align*}
T_{m+1,k}(\rho )&=\sum_{i=0}^\infty b_i^{m,k-1}\rho ^{2m-(k-1)+2i+1}+\sinh \rho \sum_{i=0}^\infty b_i^{m,k}(2m-k+2i)\rho ^{2m-k+2i-1}\\
&\quad -(2m-k)\cosh \rho \sum_{i=0}^\infty b_i^{m,k}\rho ^{2m-k+2i} \\
&=\sum_{i=0}^\infty b_i^{m,k-1}\rho ^{2(m+1)-k+2i}+\sum_{i,r=0}^\infty b_i^{m,k}\Big(\frac{2m-k+2i}{(2r+1)!}-\frac{2m-k}{(2r)!}\Big)\rho ^{2m-k+2(i+r)}\\
&=\sum_{i=0}^\infty b_i^{m,k-1}\rho ^{2(m+1)-k+2i}+b_0^{m,k}(2m-k)\sum_{r=1}^\infty  \Big(\frac{1}{(2r+1)!}-\frac{1}{(2r)!}\Big)\rho ^{2m-k+2r}\\
&\quad +\sum_{i=1}^\infty\sum_{r=0}^\infty b_i^{m,k}\Big(\frac{2m-k+2i}{(2r+1)!}-\frac{2m-k}{(2r)!}\Big)\rho ^{2m-k+2(i+r)}\\
&=\sum_{i=0}^\infty b_i^{m,k-1}\rho ^{2(m+1)-k+2i}+b_0^{m,k}(2m-k)\sum_{r=0}^\infty  \Big(\frac{1}{(2r+3)!}-\frac{1}{(2r+2)!}\Big)\rho ^{2m-k+2r+2}\\
&\quad +\sum_{i=0}^\infty\sum_{r=0}^\infty b_{i+1}^{m,k}\Big(\frac{2m-k+2i+2}{(2r+1)!}-\frac{2m-k}{(2r)!}\Big)\rho ^{2m-k+2(i+r)+2}\\
&=\sum_{i=0}^\infty b_i^{m,k-1}\rho ^{2(m+1)-k+2i}+b_0^{m,k}(2m-k)\sum_{r=0}^\infty  \Big(\frac{1}{(2r+3)!}-\frac{1}{(2r+2)!}\Big)\rho ^{2(m+1)-k+2r}\\
&\quad +\sum_{i=0}^\infty\sum_{r=0}^ib_{i-r+1}^{m,k}\Big(\frac{2m-k+2(i-r)+2}{(2r+1)!}-\frac{2m-k}{(2r)!}\Big)\rho ^{2(m+1)-k+2i}\\
&=\sum_{i=0}^\infty b_i^{m+1,k}\rho ^{2(m+1)-k+2i},\quad \rho >0,
\end{align*}
for certain $(b_i^{m+1,k})_{i=0}^\infty\subseteq \mathbb R$, for which  the series is absolutely convergent, for every $\rho >0$. Hence the proof of $(i)$ is finished.

To show $(ii)$ we fix $m\in \mathbb N$, $k=1,\ldots, m$, and $j=0,\ldots, k$, and  write
$$
Q_j^{m,k}(\rho )=\sum_{r=0}^{m-k}a_r^{m,k,j}(\sinh \rho )^{m-k-r}(\cosh \rho )^{r},\quad \rho >0,
$$
with $a_r^{m,k,j}\in \mathbb R$, $r=0,\ldots, m-k$. We observe that, for each $r=0,\ldots, m-k$,
$$
\frac{(\sinh \rho )^{m-k-r}(\cosh \rho )^r}{(\sinh \rho )^{2m-k}}=\frac{(\cosh \rho )^r}{(\sinh \rho )^{m+r}}=2^me^{-m\rho }\frac{(1+e^{-2\rho })^r}{(1-e^{-2\rho })^{m+r}},\quad \rho >0.
$$
Then, if $c_j^{m,k}=2^m\sum_{r=0}^{m-k}a_r^{m,k,j}$ we have that
$$
\frac{Q_j^{m,k}(\rho )}{(\sinh \rho )^{2m-k}}-c_j^{m,k}e^{-m\rho }=2^me^{-m\rho }\sum_{r=0}^{m-k}a_r^{m,k,j}\Big(\frac{(1+e^{-2\rho })^r}{(1-e^{-2\rho })^{m+r}}-1\Big),\quad \rho >0,
$$
and thus, when $\rho \geq 1$,
\begin{align*}
 \Big|\frac{Q_j^{m,k}(\rho )}{(\sinh \rho )^{2m-k}}-c_j^{m,k}e^{-m\rho }\Big|&\leq Ce^{-m\rho }\sum_{r=0}^{m-k}\frac{(1+e^{-2\rho })^r-(1-e^{-2\rho })^{m+r}}{(1-e^{-2\rho })^{m+r}}\\
 &\hspace{-3cm}\leq  Ce^{-m\rho }\sum_{r=0}^{m-k}\big((1+e^{-2\rho })^r-(1-e^{-2\rho })^{m+r}\big)\\
  &\hspace{-3cm}=Ce^{-m\rho }\sum_{r=0}^{m-k}\Big(\sum_{\ell =1}^r\binom{r}{\ell}e^{-2\ell\rho}-\sum_{\ell =1}^{m+r}(-1)^\ell \binom{m+r}{\ell}e^{-2\ell\rho }\Big)\leq Ce^{-(m+2)\rho },
\end{align*}
for certain $C>0$ that can depend on $m$ but not on $j$ or $k$. Thus,
we conclude that there exists $C>0$ independent of $j$ and $k$ such that
$$
 \Big|\frac{Q_j^{m,k}(\rho )}{(\sinh \rho )^{2m-k}}-c_j^{m,k}e^{-m\rho }\Big|\leq Ce^{-(m+2)\rho },\quad \rho \geq 1.
$$
Note that from this estimation property in $(iii)$ follows easily. 

Finally, let us show that, for every $m\in \mathbb N$, and $k=1,\ldots, m$,
\begin{equation}\label{cmk}
|c_k^{m,k}|=(-1)^{m+k}c_k^{m,k}.
\end{equation}
Recall that $c^{m,k}_k=2^m\sum_{r=0}^{m-k}a_r^{m,k,k}$. We proceed again by induction on $m$.  

Since $T_{1,1}(\rho )=\rho $, we have that $Q_1^{1,1}(\rho )=1$. Thus, $c_1^{1,1}=2$, so the property is verified for $m=1$. Moreover, it is obvious from \eqref{Tm} that, for each $m\in \mathbb N$, $Q^{m,m}_m(\rho )=1$, $\rho >0$. Hence, $c_m^{m,m}=2^m$ and \eqref{cmk} is trivially satisfied.

Suppose now that \eqref{cmk} holds for certain $m\in \mathbb N$, and each $k=1,\ldots, m$. From \eqref{Q11} we get
\begin{align*}
Q_1^{m+1,1}(\rho )&=\sum_{r=0}^{m-2}a_r^{m,1,1}(m-1-r)(\sinh \rho )^{m-1-r}(\cosh \rho )^{r+1}\\
&\quad +\sum_{r=1}^{m-1}ra_r^{m,1,1}(\sinh \rho )^{m+1-r}(\cosh \rho )^{r-1}\\
&\quad -(2m-1)\sum_{r=0}^{m-1}a_r^{m,1,1}(\sinh \rho )^{m-1-r}(\cosh \rho )^{r+1},\quad \rho >0.
\end{align*}
Thus, 
\begin{align*}
    2^{-(m+1)}c_1^{m+1,1}&=\sum_{r=0}^{m-2}a_r^{m,1,1}(m-1-r)+\sum_{r=1}^{m-1}ra_r^{m,1,1}-(2m-1)\sum_{r=0}^{m-1}a_r^{m,1,1}\\
    &=-m\sum_{r=0}^{m-1}a_r^{m,1,1}=-2^{-m}mc_1^{m,1},
\end{align*}
and it follows that $|c_1^{m+1,1}|=2m|c_1^{m,1}|=(-1)^{m+1}2mc_1^{m,1}=(-1)^{m+2}c_1^{m+1,1}$.

On the other hand, according to \eqref{Qkk}, for $k=2,\ldots, m$, we have that
\begin{align*}
Q_k^{m+1,k}(\rho )&=\sum_{r=0}^{m-k+1}a_r^{m,k-1,k-1}(\sinh \rho )^{m-k+1-r}(\cosh \rho )^r\\
&\quad +\sum_{r=0}^{m-k-1}a_r^{m,k,k}(m-k-r)(\sinh \rho )^{m-k-r}(\cosh \rho )^{r+1}\\
&\quad +\sum_{r=1}^{m-k}ra_r^{m,k,k}(\sinh \rho )^{m-k-r+2}(\cosh \rho )^{r-1}\\
&\quad -(2m-k)\sum_{r=0}^{m-k}a_r^{m,k,k}(\sinh \rho )^{m-k-r}(\cosh \rho )^{r+1},\quad \rho >0.
\end{align*}
Then, for every $k=2,\ldots, m$,
\begin{align*}
    2^{-(m+1)}c_k^{m+1,k}&\\
    &\hspace{-1cm}=\sum_{r=0}^{m-k+1}a_r^{m,k-1,k-1}+\sum_{r=0}^{m-k-1}a_r^{m,k,k}(m-k-r)+\sum_{r=1}^{m-k}ra_r^{m,k,k}-(2m-k)\sum_{r=0}^{m-k}a_r^{m,k,k}\\
    &\hspace{-1cm}=2^{-m}c_{k-1}^{m,k-1}-m\sum_{r=0}^{m-k}a_r^{m,k,k}=2^{-m}c_{k-1}^{m,k-1}-2^{-m}mc_k^{m,k}\\\
    &\hspace{-1cm}=(-1)^{m+k+1}2^{-m}(|c_{k-1}^{m,k-1}|+m|c_k^{m,k}|),
\end{align*}
that is, $|c_k^{m+1,k}|=2(|c_k^{m,k-1}|+m|c_k^{m,k}|)=(-1)^{m+k+1}c_k^{m+1,k}$. 
\end{proof}

\end{document}
